% Part: computability
% Chapter: computability-theory
% Section: reducibility

\documentclass[../../../include/open-logic-section]{subfiles}

\begin{document}

\olfileid[pt]{cmp}{thy}{red}
\olsection{Redutibilidade}

\begin{explain}
Sabemos agora que existe pelo menos um conjunto, $K_0$, que é
computavelmente enumerável, mas não computável. Deve estar claro que
há outros. O método da redutibilidade fornece um recurso poderoso para
mostrar que outros conjuntos têm essas propriedades, sem que seja
preciso retornar constantemente aos primeiros princípios.

Em termos gerais, uma ``redução'' de um conjunto $A$ a um conjunto~$B$
é um método para transformar respostas sobre se !!{element}s pertencem
ou não a~$B$ em respostas sobre se !!{element}s pertencem ou não
a~$A$. Vamos nos concentrar em uma noção chamada ``redutibilidade
muitos-para-um'', mas há muitas outras noções de redutibilidade, com
propriedades variadas. As noções de redutibilidade também são centrais
para o estudo da complexidade computacional, em que questões de
eficiência também precisam ser consideradas. Por exemplo, diz-se que
um conjunto é ``NP-completo'' se ele pertence a NP e todo problema em NP
pode ser reduzido a ele, usando uma noção de redução semelhante à
descrita abaixo, apenas com a exigência adicional de que a redução
possa ser calculada em tempo polinomial.

Já usamos implicitamente a noção de redução. Defina o conjunto~$K$ por
\[
K = \Setabs{x}{\cfind{x}(x) \fdefined},
\]
O.s.\@, $K = \Setabs{x}{x \in W_x}$. Nossa demonstração de que o
problema da parada é insolúvel (\olref[hlt]{thm:halting-problem})
mostra de maneira mais direta que $K$ não é computável. Lembre-se de
que $K_0$ é o conjunto
\[
K_0 = \Setabs{\tuple{e, x}}{\cfind{e}(x) \fdefined },
\]
O.s.\@, $K_0 = \Setabs{\tuple{x,e}}{x \in W_e}$. É fácil estender
qualquer demonstração da não computabilidade de~$K$ a uma demonstração
da não computabilidade de~$K_0$: se $K_0$ fosse computável, poderíamos
decidir se !!a{element}~$x$ pertence ou não a $K$ simplesmente
perguntando se o par $\tuple{x, x}$ pertence ou não a $K_0$. A
função~$f$ que associa $x$ a $\tuple{x, x}$ é um exemplo de
\emph{redução} de $K$ a $K_0$.
\end{explain}

\begin{defn}
Sejam $A$ e $B$ conjuntos de números naturais. Uma função computável
$f\colon \Nat \to \Nat$ é uma \emph{redução muitos-para-um} de $A$
a~$B$ se, e somente se, para todo número natural~$x$,
\[
x \in A \quad \text{se, e somente se,} \quad f(x) \in B.
\]
Se existe tal redução $f$, dizemos que $A$ é \emph{redutível
muitos-para-um} a~$B$, e escrevemos $A \leq_m B$. Se $A$ é redutível
muitos-para-um a $B$ e vice-versa, dizemos que $A$ e $B$ são
\emph{equivalentes muitos-para-um}, e escrevemos $A \equiv_m B$.
\end{defn}

Se a função $f$ da definição acima for injetora, diz-se que $A$~é
\emph{redutível um-para-um} a~$B$. A maioria das reduções descritas
abaixo satisfaz essa exigência mais forte, mas não usaremos esse fato.

\begin{digress}
É verdade, embora de modo algum óbvio, que a redutibilidade um-para-um
é realmente uma exigência mais forte do que a redutibilidade
muitos-para-um. Em outras palavras, existem conjuntos infinitos $A$
e~$B$ tais que $A$ é redutível muitos-para-um a~$B$, mas não é
redutível um-para-um a~$B$.
\end{digress}

\end{document}

